\begin{derivation}[从\eqnrefs{eq:sm-higgs-vector-rules-dp43}到\eqnrefs{eq:higgs-to-vv-rates-dp11}]
正文\eqnrefs{eq:sm-higgs-vector-rules-dp43} 给出 $hWW$ 与 $hZZ$ 顶角。对两个有质量矢量玻色子的三个物理极化求和，利用在壳二体运动学整理纵向贡献，再乘二体相空间并对 $ZZ$ 加入相同粒子因子，即得到正文\eqnrefs{eq:higgs-to-vv-rates-dp11}。

记 $V=W$ 或 $Z$，并定义
\begin{equation*}
M_V\equiv m_Vc^2,
\qquad
x_V\equiv\frac{m_V^2}{m_h^2}.
\end{equation*}
质量本征基的 Higgs--矢量顶角具有
\begin{equation*}
\ii g_{hVV}\eta_{\mu\nu},
\qquad
g_{hVV}=\frac{2M_V^2}{v_E}.
\end{equation*}
因此固定两条末态极化时
\begin{equation*}
\mathcal M_{\lambda_1\lambda_2}
=g_{hVV}\,
\epsilon_{1}^{(\lambda_1)*}\!\cdot
\epsilon_{2}^{(\lambda_2)*}.
\end{equation*}
每个在壳有质量矢量具有三个物理极化，完备关系为
\begin{equation*}
P_{\mu\nu}(k)
\equiv
\sum_\lambda
\epsilon_\mu^{(\lambda)}(k)
\epsilon_\nu^{(\lambda)*}(k)
=-\eta_{\mu\nu}
+\frac{k_\mu k_\nu}{m_V^2c^2}.
\end{equation*}
于是极化求和后的矩阵元平方为
\begin{align*}
\sum_{\lambda_1,\lambda_2}|\mathcal M|^2
&=g_{hVV}^2
\eta^{\mu\nu}\eta^{\rho\sigma}
P^{(1)}_{\mu\rho}P^{(2)}_{\nu\sigma}\\
&=g_{hVV}^2
\left[
4-\frac{k_1^2}{m_V^2c^2}
-\frac{k_2^2}{m_V^2c^2}
+\frac{(k_1\!\cdot k_2)^2}{m_V^4c^4}
\right].
\end{align*}
两条末态均在壳，故 $k_1^2=k_2^2=m_V^2c^2$，从而
\begin{equation*}
\sum|\mathcal M|^2
=g_{hVV}^2
\left[2+\frac{(k_1\!\cdot k_2)^2}{m_V^4c^4}\right].
\end{equation*}
Higgs 质量壳满足 $(k_1+k_2)^2=m_h^2c^2$，因此
\begin{align*}
2k_1\!\cdot k_2
&=m_h^2c^2-k_1^2-k_2^2\\
&=(m_h^2-2m_V^2)c^2,
\end{align*}
即
\begin{equation*}
\frac{(k_1\!\cdot k_2)^2}{m_V^4c^4}
=\frac{(1-2x_V)^2}{4x_V^2}.
\end{equation*}
代回后
\begin{align*}
2+\frac{(1-2x_V)^2}{4x_V^2}
&=\frac{8x_V^2+1-4x_V+4x_V^2}{4x_V^2}\\
&=\frac{1-4x_V+12x_V^2}{4x_V^2}.
\end{align*}
再乘 $g_{hVV}^2=4M_V^4/v_E^2$，得到熟悉的运动学多项式
\begin{equation*}
\sum|\mathcal M|^2
=\frac{(m_hc^2)^4}{v_E^2}
\left(1-4x_V+12x_V^2\right).
\end{equation*}
二体相空间的末态速度因子为
\begin{equation*}
\beta_V=\sqrt{1-4x_V}.
\end{equation*}
把该矩阵元平方、$\beta_V$ 和 $1/v_E^2=\sqrt2G_F^{(E)}$ 代入一体衰变母式，得到\eqnrefs{eq:higgs-to-vv-rates-dp11}。$ZZ$ 末态的两粒子完全相同，因此其相空间要除以 $2!$；$W^+W^-$ 可区分，不含这一因子，所以两个宽度的整体前因子相差二倍。
\end{derivation}
